\chapter{Methodology and Experimental Design}\label{ch:methodology}

This chapter details the experimental framework used to evaluate the influence of embedding-model capacity on retrieval performance. The methodology follows a hypothesis-driven, research-TDD approach: establishing a baseline pipeline, validating it through pilot testing, and systematically expanding experimental conditions to test the core hypothesis.

\section{Dataset Description and Preprocessing}

\subsection{The iPalpiti Classical Archive}
The primary dataset consists of the iPalpiti classical music archive, a curated collection of live orchestral and chamber music performances. This single-domain archive features exclusively classical works performed by the iPalpiti festival orchestra. The collection comprises a total of [TBD: Insert Final Track Count] tracks, representing approximately [TBD: Insert Total Hours] hours of audio recordings.

\subsection{Metadata and Label Taxonomy}
Ground truth for proxy retrieval tasks is derived from editorial metadata associated with each recording. The composer metadata serves as an objective label for sanity-check retrieval tasks, allowing us to verify if models can cluster tracks by composition. For the primary musical proxy task, tracks are annotated with abstract character labels (e.g., ``Heavy'' versus ``Light''), which capture the timbral and structural qualities relevant to the archive's curatorial goals.

\subsection{Audio Preprocessing Pipeline}
To ensure internal validity, the audio preprocessing pipeline is held constant across all embedding models. All audio files are first resampled to [TBD: e.g., 16kHz/32kHz] to satisfy the input requirements of the selected models. Following resampling, the channel format is standardized to [TBD: Mono/Stereo], and [TBD: Description of normalization, e.g., peak normalization] is applied to ensure consistent amplitude levels. Finally, the audio is segmented at the track level; for models requiring fixed-length inputs, [TBD: Describe truncation or windowing strategy] is employed to extract representative segments.

\section{Selection of Embedding Models (Independent Variable)}

Models are selected to represent discrete points along the capacity spectrum, ranging from compact convolutional architectures to large-scale transformer models. This selection allows us to treat model capacity as the primary independent variable.

\subsection{Experimental Matrix}
The experiment compares performance across three distinct model families representing low, medium, and high capacity tiers. Table \ref{tab:model_matrix} summarizes the selected architectures and their operational status within the research timeline.

\begin{table}[h]
\centering
\begin{tabular}{|l|l|l|l|}
\hline
Model Family & Representative Model & Capacity ($\approx |\theta|$) & Status \\ \hline
Compact / CNN & [TBD: e.g., VGGish / MusiCNN] & Low ($< 10M$) & Planned \\ \hline
Transformer (Base) & [TBD: e.g., AST / MERT-95M] & Medium ($\approx 80-100M$) & Planned \\ \hline
Transformer (Large) & MERT-330M & High ($> 300M$) & Pilot Executed \\ \hline
\end{tabular}
\caption{Model Selection Matrix representing the independent variable (Capacity).}
\label{tab:model_matrix}
\end{table}

\subsection{Low-Capacity Baselines}
To establish a performance floor, we utilize compact architectures typically optimized for mobile or real-time inference. The representative model for this category is [Model A (TBD)], a [Description of architecture, e.g., CNN-based] architecture trained using supervised objectives.

\subsection{High-Capacity Transformers}
To evaluate the relationship between model capacity and retrieval performance, we include large-scale transformer architectures. The primary model selected for this category is the Music Understanding Model with Large-Scale Self-Supervised Training (MERT), a transformer-based architecture pretrained using masked language modeling objectives on musical audio. Additionally, [TBD: Additional Large Model] is included to provide a comparative reference within the high-capacity tier.

\section{The Retrieval Pipeline}

\subsection{Inference Engine}
The retrieval system is implemented as a modular backend pipeline that isolates the embedding extraction process from downstream retrieval tasks. The system is built using the PyTorch and HuggingFace Transformers frameworks and executed on [TBD: Hardware Specs, e.g., NVIDIA A100 / RTX 3090] hardware to ensure efficient batch processing.

\subsection{Similarity and Indexing}
Similarity between item embeddings is computed using the Cosine Similarity metric. Given the small scale of the target dataset, we employ exact nearest-neighbor search rather than approximate indexing methods. This ensures that retrieval results are deterministic and free from approximation errors.

\section{Pilot Validation: Pipeline Unit Test}
Before full-scale execution, a pilot experiment was conducted to validate the integrity of the retrieval pipeline. 
This corresponds to the validation phase of the research-oriented test-driven development framework.


\subsection{Pilot Configuration}
The pilot experiment utilized a subset of 10 randomly selected tracks from the iPalpiti archive. The high-capacity MERT-330M model was employed to verify that the pipeline could handle complex transformer embeddings. The primary objective was to confirm that the pipeline produces consistent embeddings and correctly computes pairwise similarity and ranking.

\subsection{Pilot Success Criteria}
The pilot was evaluated against three specific criteria. First, the system must successfully generate embeddings for all input tracks without errors. Second, the retrieval results must be deterministic and reproducible across multiple runs. Third, the distribution of cosine similarity scores must exhibit variance, ensuring that the model has not suffered from embedding collapse where all vectors converge to a single point.

\section{Evaluation Protocol (Dependent Variables)}

\subsection{Relevance Judgments (Ground Truth)}
Relevance is determined by metadata matching as defined in Chapter 3. The first proxy task, the Sanity Proxy, defines relevance based on shared composer identity; a retrieved item is considered relevant only if it shares the same composer as the query item. The second task, the Primary Musical Proxy, defines relevance based on abstract musical character labels; items are relevant if they share the same descriptive character (e.g., ``Heavy'') as the query.

\subsection{Performance Metrics}
Ranking quality is measured using standard Information Retrieval metrics. We report Precision@K to measure the proportion of relevant items in the top $K$ results, and Recall@K to measure the proportion of total relevant items retrieved. Additionally, Normalized Discounted Cumulative Gain (NDCG@K) is used to evaluate ranking quality, penalizing relevant items that appear lower in the retrieved list.

\section{Experimental Phases}

The study proceeds in three phases. Phase I, the Sanity Check, validates whether models can cluster tracks by objective, high-level metadata such as composer identity. Phase II, the Musical Character Retrieval task, tests the core hypothesis by evaluating retrieval performance on abstract musical qualities. Finally, Phase III performs a Capacity-Performance Analysis, correlating the measured retrieval metrics with model capacity to determine if performance gains saturate or regress at higher capacities within the small-scale archival context.
